\documentclass[12pt,a4paper]{amsart}

%% ── Packages ─────────────────────────────────────────────────────────────────
\usepackage{amsmath,amssymb,amsthm}
\usepackage{mathtools}
\usepackage{hyperref}
\usepackage{cleveref}
\usepackage{enumitem}
\usepackage[margin=2.5cm]{geometry}

%% ── Theorem environments ──────────────────────────────────────────────────────
\theoremstyle{plain}
\newtheorem{theorem}{Theorem}[section]
\newtheorem{lemma}[theorem]{Lemma}
\newtheorem{corollary}[theorem]{Corollary}
\newtheorem{proposition}[theorem]{Proposition}

\theoremstyle{definition}
\newtheorem{definition}[theorem]{Definition}
\newtheorem{remark}[theorem]{Remark}

%% ── Macros ────────────────────────────────────────────────────────────────────
\newcommand{\R}{\mathbb{R}}
\newcommand{\Q}{\mathbb{Q}}
\newcommand{\Z}{\mathbb{Z}}
\newcommand{\N}{\mathbb{N}}
\newcommand{\C}{\mathbb{C}}

% Weil form and spaces
\newcommand{\HL}{L^2(I_L)}
\newcommand{\IL}{I_L}
\newcommand{\QW}{Q_W}
\newcommand{\qbar}{\bar{q}}
\newcommand{\qtilde}{\tilde{q}}

% Operators
\newcommand{\SbL}{S_{b,L}}
\newcommand{\CbL}{C_{b,L}}
\newcommand{\Ctau}{C_{\tau,1}}

% Constants
\newcommand{\ctwo}{c_2}              % = log2/sqrt(2)
\newcommand{\kappae}{\kappa_{\mathrm{edge}}}
\newcommand{\Lzero}{L_0}             % = 2^{-30}

% Intervals
\newcommand{\Iminus}{I_-}
\newcommand{\Izero}{I_0}
\newcommand{\Iplus}{I_+}
\newcommand{\Eminus}{E_-}
\newcommand{\Eplus}{E_+}

% Norms
\newcommand{\norm}[1]{\left\|#1\right\|}
\newcommand{\ip}[2]{\left\langle #1,\, #2\right\rangle}

% Legendre
\newcommand{\Pn}{P_n}
\newcommand{\Jij}{J_{ij}(\tau)}
\newcommand{\Eij}{E_{ij}(\tau)}

%% ── Title ─────────────────────────────────────────────────────────────────────
\title{Structural Lemmas for the First-Prime Window\\
       of the Weil Quadratic Form}

\author{Lin Tao}

\date{\today}

\begin{document}

\maketitle

\begin{abstract}
We prove six structural results for the local Weil quadratic form
$Q_W^L$ on $L^2(-L,L)$ in the interval
$\tfrac{1}{2}\log 2 < L < \tfrac{1}{2}\log 3$,
the \emph{first-prime window} where only the prime $n = 2$ contributes
to the explicit formula.
Our main technical contribution (\cref{thm:prime-legendre}) is an
algebraisation showing that the matrices $J_{ij}(\tau)$ and $E_{ij}(\tau)$
encoding the coupling between the prime-2 layer and the Legendre basis lie in
$\mathbb{Q}[\tau]$ and are computable without numerical quadrature.
We also prove a pure-rational absorption certificate (\cref{thm:rational-cert}):
$V + P_{2,7/20} \ge \tfrac{69}{100}V \ge 0$, with the key step
$87^{16}\cdot 68^5 < 1701^5\cdot 32^{16}$ verified in Lean~4/Mathlib.
We give an exact spectral description of $C_{b,L}$ (\cref{thm:overlap}),
falsify Path~A by explicit certified negative witnesses (\cref{thm:path-a-neg}),
and provide a conditional Path~B Schur criterion (\cref{thm:schur}):
if $b_L > 0$ and $b_L F - R_\eta \succ 0$, then $\lambda(7/20) > 0$;
the verification of these matrix conditions is ongoing work.
All results are independent of the main conjecture FP-0.35
($\lambda(7/20) > 0$) and do not imply the Riemann Hypothesis.
\end{abstract}

\tableofcontents

%% ══════════════════════════════════════════════════════════════════════════════
\section{Introduction}
\label{sec:intro}
%% ══════════════════════════════════════════════════════════════════════════════

\subsection{Background}

The Weil explicit formula associates to any even Schwartz function $f$ a
functional
\[
  W(f) = \hat f(0)\log\pi - \sum_\rho f(\gamma_\rho)
         - \sum_{n \ge 2} \frac{\Lambda(n)}{\sqrt n}
           \bigl(f(\log n) + f(-\log n)\bigr),
\]
where the sum over $\rho$ runs over the non-trivial zeros $\rho = \tfrac12 +
i\gamma_\rho$ of the Riemann zeta function.
A classical result due to Weil~\cite{Weil1952} states that the
\emph{Riemann Hypothesis} (RH) is equivalent to $W(f * \widetilde f) \ge 0$
for all admissible $f$, where $\widetilde f(x) = \overline{f(-x)}$.

The \emph{local Weil quadratic form} restricts attention to test functions
supported in $(-L,L)$:
\[
  Q_W^L(v) = W(v * \widetilde v), \qquad v \in L^2(-L,L).
\]
Yoshida~\cite{Yoshida1992} and Bombieri~\cite{Bombieri2000} proved that
$Q_W^L \ge 0$ for all $L$ in a certain small-support regime.
Suzuki~\cite{Suzuki2026} gave a concrete Galerkin realization and
developed the screw-function framework that we adopt;
Groskin~\cite{Groskin2026} provided a finite Guinand--Weil dictionary
and archimedean tail estimates closely related to the present work;
and subsequent numerical computation~\cite{Kim2026} established positivity up to
$L \approx 0.35$ by machine computation.

\subsection{Scope and independent value of the results}

We work entirely in the \emph{first-prime window}
\[
  \tfrac{1}{2}\log 2 < L < \tfrac{1}{2}\log 3,
\]
where the prime-side contribution to $Q_W^L$ consists of the single term
$n = 2$.
Using Suzuki's conventions~\cite{Suzuki2026}, the $n$-th prime term is the
quadratic form of the truncated shift operator
\[
  Q_{n,L}(v) = -\frac{\Lambda(n)}{\sqrt n}\ip{C_{\log n,L}\,v}{v},
\]
where $C_{b,L} = S_{b,L} + S_{b,L}^*$ and
$(S_{b,L}v)(x) = \mathbf{1}_{I_L}(x)\,v(x-b)$.

\medskip
\textbf{Independent value of the results.}
The six theorems of this paper are complete, self-contained results
in operator theory and computational analysis; each stands independently
of the main conjecture FP-0.35.
In particular:
\begin{itemize}[leftmargin=2em]
\item \emph{Theorem~\ref{thm:prime-legendre}}
  (first-prime Legendre matrix algebraisation) is the technical core of this
  paper. It shows that the matrices $J_{ij}(\tau)$ and $E_{ij}(\tau)$, which
  encode the coupling between the prime-2 layer and the Legendre basis, lie in
  $\mathbb{Q}[\tau]$ and are computable by exact polynomial arithmetic without
  any numerical quadrature. This makes the prime layer algebraically tractable
  and underlies the Path~B Schur criterion (Theorem~\ref{thm:schur}).
\item \emph{Theorem~\ref{thm:rational-cert}} (pure-rational certificate)
  provides a pure-rational certificate for $V + P_{2,7/20} \ge (69/100)V$
  whose key steps are machine-verified in Lean~4 via Mathlib. Specifically,
  the five integer comparisons (including $87^{16} \cdot 68^5 < 1701^5 \cdot 32^{16}$)
  are verified by \texttt{native\_decide}, and the transcendental bounds
  $\log 2 < 7/10$ and $\sqrt{2} > 7/5$ are verified using
  \texttt{Real.sum\_le\_exp\_of\_nonneg} and \texttt{Real.sqrt\_lt\_sqrt}
  from Mathlib's analysis library. The algebraic reductions connecting these
  bounds are carried out by hand in the proof.
\item \emph{Theorem~\ref{thm:overlap}} (spectral decomposition of $C_{b,L}$)
  gives a complete, elementary description of the first-prime layer as an
  operator: its spectrum is $\{-1,0,1\}$ with all eigenvalues of infinite
  multiplicity, and it is not compact. This rules out several naive reduction
  strategies and is used in the proof of Theorem~\ref{thm:prime-legendre}.
\item \emph{Theorem~\ref{thm:path-a-neg}} (Path~A falsification)
  proves that the potential redistribution strategy fails at $L = 7/20$ with
  explicit interval-certified negative witnesses, showing the absorption
  threshold is sharp.
\end{itemize}

\medskip
\textbf{Scope boundary.}
\emph{The results of this paper do not imply the Riemann Hypothesis.}
Even if the main conjecture FP-0.35 ($\lambda(7/20) > 0$) is established,
the passage from finite-scale positivity to global positivity (which would
imply RH) requires techniques not developed here.
The paper is not incomplete for failing to resolve FP-0.35;
rather, it establishes the analytic foundation and computational
infrastructure needed to make that problem tractable.

\subsection{Organisation}

\Cref{sec:prelim} fixes notation.
\Cref{sec:overlap} establishes the exact spectral decomposition of $C_{b,L}$
(Theorems~\ref{thm:overlap}--\ref{thm:no-small-pert}).
\Cref{sec:absorption} gives the endpoint potential absorption estimate
(Theorem~\ref{thm:absorption}) and its pure-rational certificate
(Theorem~\ref{thm:rational-cert}).
\Cref{sec:path-b} develops the Path~B algebraic machinery
(Theorems~\ref{thm:prime-legendre}--\ref{thm:schur}).
\Cref{sec:falsification} proves that Path~A is strictly falsified
(Theorem~\ref{thm:path-a-neg}).
\Cref{sec:conjecture} states the main conjecture FP-0.35 and the remaining
obligations.

%% ══════════════════════════════════════════════════════════════════════════════
\section{Preliminaries}
\label{sec:prelim}
%% ══════════════════════════════════════════════════════════════════════════════

We write $I_L = (-L,L)$ and $H_L = L^2(I_L)$ with the standard inner product
$\ip{f}{g} = \int_{-L}^L f(x)\overline{g(x)}\,dx$.
Functions $v \in H_L$ are zero-extended to $\R$.

The \emph{log-endpoint potential} is
$V(x) = -\tfrac{1}{2}\log(1-x^2)$ on $(-1,1)$ (in the scaled model).
After the scaling $w(t) = v(Lt)$ and $\tau = \log 2/L$, the first-prime
operator becomes $\mathcal P_{2,L} = -\ctwo C_{\tau,1}$ where
$\ctwo = \log 2/\sqrt 2$ and $\Ctau$ is the truncated shift on $L^2(-1,1)$.

The \emph{Archimedean kernel} $K_L$ is the integral operator with kernel
$k_L(x,y) = -L\,r''(L(x-y))$, where $r''$ is the second derivative of
Riemann's $r$-function.
The \emph{closed quadratic form} is
\[
  \qbar_L(v) = \ip{Tv}{v} + \ip{Vv}{v} + \ip{K_Lv}{v}
               + \ip{\mathcal P_{2,L}v}{v} - c_L\norm{v}^2,
\]
where $T$ is the kinetic term (harmonic-number diagonal in Legendre basis)
and $c_L$ is the Weil constant at radius $L$.

%% ══════════════════════════════════════════════════════════════════════════════
\section{Spectral Decomposition of the Truncated Shift}
\label{sec:overlap}
%% ══════════════════════════════════════════════════════════════════════════════

\begin{theorem}[Single-step overlap decomposition]
\label{thm:overlap}
Let $b > 0$ and $b/2 < L < b$.
Define
\[
  \Iminus = (-L,\, L-b),\quad
  \Izero  = (L-b,\, -L+b),\quad
  \Iplus  = (-L+b,\, L).
\]
Then $H_L = L^2(\Iminus) \oplus L^2(\Izero) \oplus L^2(\Iplus)$,
the map $T_b\colon L^2(\Iminus) \to L^2(\Iplus)$, $(T_bh)(x) = h(x-b)$,
is unitary, and in this decomposition with the identification $T_b$,
\[
  \CbL \simeq
  \begin{pmatrix} 0 & 1 \\ 1 & 0 \end{pmatrix}
  \otimes I_{L^2(\Iminus)}
  \;\oplus\; 0_{L^2(\Izero)}.
\]
In particular $\sigma(\CbL) = \{-1,0,1\}$, $\norm{\CbL} = 1$,
all eigenvalues have infinite multiplicity, and $\CbL$ is neither compact
nor of finite rank.
\end{theorem}

\begin{proof}
Since $b/2 < L < b$, all three intervals have positive length and no three
points $x, x+b, x+2b$ lie simultaneously in $I_L$.
For $h \in L^2(\Iminus)$ set $h_+ = T_bh \in L^2(\Iplus)$; then
$\CbL(h \pm h_+) = \pm(h \pm h_+)$, giving eigenvectors for $\pm 1$.
Functions supported in $\Izero$ satisfy $S_{b,L}f = S_{b,L}^*f = 0$,
so $\Izero$ contributes eigenvalue~$0$.
Since $L^2(\Iminus)$ is infinite-dimensional (it is isomorphic to $L^2$ of a
non-degenerate interval), every eigenvalue $\pm 1$ has infinite multiplicity;
similarly the eigenspace for $0$ contains $L^2(\Izero)$, also infinite-dimensional.
The spectrum $\sigma(\CbL) = \{-1,0,1\}$ follows because $\CbL^2 =
\mathbf{1}_{L^2(\Iminus\cup\Iplus)}$ (the identity on the two boundary strips),
so every $\lambda\in\sigma(\CbL)$ satisfies $\lambda^2\le 1$, giving
$\sigma(\CbL)\subseteq[-1,1]$; together with the explicit eigenvalues this yields
equality.
Non-compactness: a compact operator on an infinite-dimensional Hilbert space
cannot have $\pm 1$ as eigenvalues of infinite multiplicity (its spectrum
accumulates only at~$0$); since $\CbL$ has $\pm 1$ as eigenvalues of infinite
multiplicity, it is not compact.
Finite rank is ruled out by the same argument.
\end{proof}

\begin{corollary}[Exact sign structure of the first-prime layer]
\label{cor:first-prime}
Let $b = \log 2$.
In the first-prime window $\tfrac{1}{2}\log 2 < L < \tfrac{1}{2}\log 3$,
\[
  \sigma\!\left(-\ctwo C_{\log 2,L}\right)
  = \left\{-\ctwo,\; 0,\; +\ctwo\right\},
  \qquad \ctwo = \frac{\log 2}{\sqrt 2}.
\]
Both the negative and positive spectral subspaces are infinite-dimensional.
\end{corollary}

\begin{corollary}[No $L^2$ operator-norm small perturbation at threshold]
\label{thm:no-small-pert}
For any $L \in (\tfrac{1}{2}\log 2,\, \log 2)$,
$\norm{C_{\log 2,L}} = 1$
regardless of the support length $2L - \log 2 \to 0^+$.
The operator norm jumps from $0$ to $1$ at the threshold
$L = \tfrac{1}{2}\log 2$.
\end{corollary}

%% ══════════════════════════════════════════════════════════════════════════════
\section{Endpoint Potential Absorption}
\label{sec:absorption}
%% ══════════════════════════════════════════════════════════════════════════════

\begin{theorem}[Endpoint potential absorption]
\label{thm:absorption}
Scale to $(-1,1)$: let $w(t) = v(Lt)$, $\tau = \log 2/L$,
$\varepsilon = 2 - \tau$.
For $\tfrac{1}{2}\log 2 < L < \log 2$ one has $\tau \in (1,2)$ and
$\Ctau$ couples only the boundary strips
$\Eminus = (-1,-1+\varepsilon)$ and $\Eplus = (1-\varepsilon,1)$.
Define $\kappae(L) = \tfrac{1}{2}\log\!\tfrac{1}{2\varepsilon} > 0$.
Then for every $w$ in the closed-form domain,
\begin{equation}\label{eq:absorption}
  \bigl|\ip{\Ctau w}{w}\bigr|
  \;\le\; \norm{w}_{L^2(\Eminus \cup \Eplus)}^2
  \;\le\; \kappae(L)^{-1}\ip{Vw}{w}.
\end{equation}
Consequently, if $\ctwo < \kappae(L)$,
\[
  V + \mathcal{P}_{2,L}
  \;\ge\; \left(1 - \frac{\ctwo}{\kappae(L)}\right) V \;\ge\; 0.
\]
\end{theorem}

\begin{proof}
The first inequality in~\eqref{eq:absorption} follows from the exchange
structure in \cref{thm:overlap} together with $2|ab| \le |a|^2 + |b|^2$.
For $x \in \Eminus \cup \Eplus$ one has $1 - x^2 \le 2\varepsilon$, so
$V(x) \ge \kappae(L)$.
Integrating yields the second inequality; multiplying by $\ctwo$ gives the
stated form bound.
\end{proof}

\begin{theorem}[Pure-rational absorption certificate]
\label{thm:rational-cert}
At $L = \tfrac{7}{20}$,
\[
  \frac{\ctwo}{\kappae(7/20)} < \frac{31}{100},
\]
and therefore
\[
  \boxed{V + P_{2,7/20} \;\ge\; \frac{69}{100}\,V \;\ge\; 0.}
\]
The certificate reduces to five integer comparisons, all of which
are independently machine-verified in Lean~4:
\begin{enumerate}
\item $23581 \times 10 < 34020 \times 7$
      \hfill(\,$\log 2 < 7/10$\,)
\item $34 \times 41 < 1701$
      \hfill(\,$\varepsilon < 1/41$\,)
\item $87^{16} \cdot 68^{5} < 1701^{5} \cdot 32^{16}$
      \hfill(\,$\kappae(7/20) > 8/5$\,, \textbf{key step})
\item $7^2 < 2 \times 5^2$
      \hfill(\,$\sqrt{2} > 7/5$\,)
\item $62 \times 5 \times 100 = 31 \times 8 \times 125$
      \hfill(\,ratio $= 31/100$\,)
\end{enumerate}
Each comparison is a closed proposition over $\mathbb{N}$ and is
discharged by Lean~4's \texttt{native\_decide} tactic in under one second.
\end{theorem}

\begin{proof}
\textbf{Step 1.} Using $\log 2 = 2\operatorname{arctanh}(\tfrac{1}{3})$
and the positive-term tail bound,
\begin{equation}\label{eq:log2-bounds}
  \frac{842}{1215} < \log 2 < \frac{23581}{34020} < \frac{7}{10}.
\end{equation}
Hence $\varepsilon = 2 - \tfrac{20\log 2}{7} < 2 - \tfrac{20}{7}\cdot\tfrac{842}{1215}
= \tfrac{34}{1701} < \tfrac{1}{41}$.

\textbf{Step 2.}
One verifies the integer inequality
\begin{align*}
  87^{16} \cdot 68^5
  &= \underbrace{15662194229696887109605438749172023641088}_{39\text{ digits}} \\
  &\;<\;
  \underbrace{17215562650769453014744867057217543602176}_{39\text{ digits}}
  = 1701^5 \cdot 32^{16},
\end{align*}
from which $e^{16} < (87/32)^{16} < (1701/68)^5 < (1/(2\varepsilon))^5$,
giving $\kappae(7/20) > \tfrac{8}{5}$.

\textbf{Step 3.}
From~\eqref{eq:log2-bounds} and $\sqrt{2} > \tfrac{7}{5}$,
$\ctwo < \tfrac{62}{125}$.

\textbf{Step 4.}
$\ctwo/\kappae < (62/125)/(8/5) = 31/100$.
Applying \cref{thm:absorption} completes the proof.
\end{proof}

\begin{corollary}[Potential redistribution]
\label{cor:redistribution}
The full scaled form satisfies
\[
  \qbar_{7/20}
  \;\ge\;
  \qtilde_{7/20}
  \;:=\;
  T + \tfrac{69}{100}V + K_{7/20} - c_{7/20}I
\]
in the closed-form ordering.
Proving $\qtilde_{7/20} \ge \Lzero I$ with $\Lzero > 0$ is sufficient for
FP-0.35.
\end{corollary}

%% ══════════════════════════════════════════════════════════════════════════════
\section{Path B: Prime-Layer Algebraisation and Schur Criterion}
\label{sec:path-b}
%% ══════════════════════════════════════════════════════════════════════════════

\begin{theorem}[First-prime Legendre matrix algebraisation]
\label{thm:prime-legendre}
Let $P_n$ denote the $n$-th Legendre polynomial, $\tau = \log 2/L \in (1,2)$,
and define on $L^2(-1,1)$
\[
  \Jij = \ip{\Ctau P_j}{P_i},
  \qquad
  \Eij = \ip{\Ctau P_j}{\Ctau P_i}.
\]
Then $\Jij = \Eij = 0$ when $i+j$ is odd.
When $i+j$ is even,
\begin{align}
  \Jij &= 2\int_{-1}^{1-\tau} P_i(x)\,P_j(x+\tau)\,dx, \label{eq:J}\\
  \Eij &= 2\int_{-1}^{1-\tau} P_i(x)\,P_j(x)\,dx. \label{eq:E}
\end{align}
In particular $\Jij, \Eij \in \Q[\tau]$, computable by Legendre recurrence
and exact polynomial arithmetic without any quadrature.

Sample values: $J_{00} = 4-2\tau$;
$J_{11} = \tfrac{\tau^3}{3} - 2\tau + \tfrac{4}{3}$;
$J_{02} = -\tau^3 + 3\tau^2 - 2\tau$.
\end{theorem}

\begin{proof}
The parity vanishing follows from the reflection $x \mapsto -x$ applied to
the second shift integral.
When $1 < \tau < 2$, \cref{thm:overlap} gives $\Ctau^2 = \mathbf{1}_{\Eminus\cup\Eplus}$,
so $\Eij$ is a Legendre Gram integral on two boundary strips,
yielding~\eqref{eq:E}.
Equation~\eqref{eq:J} follows by moving both shift integrals to the left
boundary strip.
The $\Q[\tau]$ membership and sample values are verified by explicit Legendre
recursion and polynomial antiderivatives.
\end{proof}

\begin{lemma}[Form boundedness of $V$ relative to $T$]
\label{lem:form-bounded}
The potential $V(x) = -\tfrac{1}{2}\log(1-x^2)$ is infinitesimally form-bounded
relative to the kinetic operator $T$ (whose Legendre basis representation has
diagonal $H_{n_j} \cdot G_{jj}$, with $H_n = \sum_{k=1}^n \tfrac{1}{k}$).
Concretely, for every $\delta > 0$ there exists $C_\delta > 0$ such that
\[
  \ip{Vv}{v} \le \delta\,\ip{Tv}{v} + C_\delta\,\|v\|^2
  \qquad \forall\, v \in D(T).
\]
Consequently the closed quadratic form $\qbar_L$ is well-defined on the
domain $D(\qbar_L) = D(T^{1/2})$, and all matrix operations in
Theorem~\ref{thm:schur} take place within this domain.
\end{lemma}

\begin{proof}
The kinetic operator $T$ acts on the Legendre basis as
$(T P_n)(x) = H_n P_n(x)$ (where $H_n = \sum_{k=1}^n k^{-1}$),
so $\ip{Tv}{v} = \sum_n H_n |\hat v_n|^2$ for $v = \sum_n \hat v_n P_n$.
Since $H_n \sim \log n$, the quadratic form of $T$ is comparable to that of
the operator $\log(1-\partial_x^2)^{1/2}$ on $(-1,1)$ with Dirichlet boundary
conditions, which in turn controls the $H^{1/2}(-1,1)$ Sobolev norm.
The logarithmic Hardy inequality on $(-1,1)$ states that for $f \in H^{1/2}_0$,
\[
  \int_{-1}^1 \frac{|f(x)|^2}{(1-x^2)\log^2 e(1-x^2)} \, dx
  \;\lesssim\; \ip{Tf}{f} + \|f\|^2.
\]
Since $V(x) = -\tfrac12\log(1-x^2) \lesssim |\log(1-x^2)|$, the potential
$\ip{Vf}{f}$ is controlled by the left side, giving the $\delta$-bound for any $\delta>0$
with an explicit $C_\delta$ depending on the Legendre harmonic numbers $H_n$.
The Archimedean kernel $K_L$ is Hilbert--Schmidt (its kernel $k_L(x,y) = -Lr''(L(x-y))$
is square-integrable on $I_L^2$), hence relatively compact
and $\delta$-form-bounded for any $\delta > 0$.
The closed form $\qbar_L$ is therefore defined on $D(T^{1/2})$ by the KLMN theorem.
\end{proof}

\begin{theorem}[Split-residual Schur criterion]
\label{thm:schur}
Write $\qbar_L = T + V + K_L + \mathcal{P}_{2,L} - c_L I$.
We denote by $\kappa_L$ a certified upper bound on the operator norm
$\|K_L\|$ in the $L^2(I_L)$ sense; such a bound is computed by the
Hilbert--Schmidt norm estimate
$\kappa_L \ge \|K_L\|_{\mathrm{HS}} \ge \|K_L\|$
using the kernel function $k_L(x,y) = -L\,r''(L(x-y))$.
Note that $\kappa_L$ is a computable rational upper bound, distinct from
$\kappae(L) = \tfrac12\log\tfrac{1}{2\varepsilon}$ which controls the
endpoint potential absorption (Theorem~\ref{thm:absorption}).
Let $\Pi_N$ be the Legendre projection onto $N$ basis functions of a fixed
parity, $\mathsf{Q}_N = I - \Pi_N$, and $d$ the first complement degree.
By \cref{thm:rational-cert}, $V + \mathcal{P}_{2,L} \ge 0$ at $L = 7/20$.
Define
\begin{gather*}
  G_{ij} = \ip{P_j}{P_i},\quad
  (T_N)_{ij} = H_{n_j}G_{ij},\\
  M^{(0)}_{ij} = \ip{(V+K_L)P_j}{P_i},\quad
  S^{(0)}_{ij} = \ip{(V+K_L)P_j}{(V+K_L)P_i},\\
  M^{(2)}_{ij} = -\ctwo J_{ij}(\tau),\quad
  S^{(2)}_{ij} = \ctwo^2 E_{ij}(\tau),\\
  R_0 = S^{(0)} - (M^{(0)})^* G^{-1} M^{(0)},\quad
  R_2 = S^{(2)} - (M^{(2)})^* G^{-1} M^{(2)},\\
  R_\eta = (1+\eta)R_0 + (1+\eta^{-1})R_2,\quad \eta \in \Q_{>0},\\
  F = T_N + M^{(0)} + M^{(2)} - (c_L + \Lzero)G,\quad
  b_L = H_d - c_L - \Lzero - \kappa_L.
\end{gather*}
If $b_L > 0$ and $b_L F - R_\eta \succ 0$, then
$\qbar_L[w] \ge \Lzero\norm{w}^2$ on the full closed-form domain
of the parity sector.
\end{theorem}

\begin{proof}
For $p \in \operatorname{Ran}\Pi_N$ and $q \in \operatorname{Ran}\mathsf{Q}_N$,
the cross-term satisfies
$\norm{\mathsf{Q}_N(V+K_L+\mathcal{P}_{2,L})p}^2 \le
 (1+\eta)\norm{\mathsf{Q}_N(V+K_L)p}^2 +
 (1+\eta^{-1})\norm{\mathsf{Q}_N\mathcal{P}_{2,L}p}^2$
by weighted Young.
The two norms have Gram matrices $R_0$ and $R_2$ respectively.
The complement lower bound $b_L > 0$ gives $\mathsf{Q}_N(\qbar_L - \Lzero)\mathsf{Q}_N \ge b_L\mathsf{Q}_N$.
Completing the square yields $(\qbar_L - \Lzero)[p+q] \ge z^*(F - b_L^{-1}R_\eta)z$,
whence $b_L F - R_\eta \succ 0$ implies the stated lower bound.
\end{proof}

%% ══════════════════════════════════════════════════════════════════════════════
\section{Strict Falsification of Path A}
\label{sec:falsification}
%% ══════════════════════════════════════════════════════════════════════════════

\begin{theorem}[Path A strict negative witnesses]
\label{thm:path-a-neg}
At $L = 7/20$, $\theta = 69/100$, $\Lzero = 2^{-30}$, the form
$\qtilde_{7/20} - \Lzero I = T + \theta V + K_{7/20} - (c_{7/20}+\Lzero)I$
satisfies
\begin{align*}
  (\qtilde_{7/20} - \Lzero I)[P_0 - P_2]
  &\in [-0.053384,\; -0.052711], \\
  (\qtilde_{7/20} - \Lzero I)[P_1 - \tfrac{1}{2}P_3]
  &\in [-0.032327,\; -0.032119].
\end{align*}
In particular, both upper endpoints are strictly negative, so the weakened
sufficient condition $\qtilde_{7/20} \ge \Lzero I$ is \emph{false}.

\begin{remark}
This does not falsify FP-0.35.
The inequality $\qbar_{7/20} \ge \qtilde_{7/20}$ holds by
\cref{cor:redistribution}, but $\qtilde_{7/20}$ having a negative direction
carries no conclusion about the sign of $\qbar_{7/20}$.
\end{remark}
\end{theorem}

\begin{proof}[Proof sketch]
The auto-correlation polynomials of $v_{\rm even} = P_0 - P_2$ and
$v_{\rm odd} = P_1 - \tfrac12 P_3$ are computed exactly from Legendre
recurrence:
\begin{align*}
  C_{v_{\rm even}}(t) &= \tfrac{12}{5} - 3t^2 + \tfrac32 t^3
                         - \tfrac{3}{40}t^5, \\
  C_{v_{\rm odd}}(t)  &= \tfrac{31}{42} - \tfrac14 t - \tfrac52 t^2
                         + \tfrac{23}{12}t^3 - \tfrac{7}{32}t^5
                         + \tfrac{5}{448}t^7.
\end{align*}
\begin{remark}[Normalisation convention for $C_v$]
\label{rem:normalisation}
We define $C_v(t) = \int_{-1}^{1-t} v(x+t)\,v(x)\,dx$, which is the
\emph{one-sided} auto-correlation (no leading factor of~$2$).
The reduction $K_L[v] = 2\int_0^2 (-L\,r''(Lt))\,C_v(t)\,dt$ uses the
explicit factor~$2$ that arises from splitting the double integral
$\iint_{I_L^2} k_L(x,y)\,v(x)v(y)\,dx\,dy$ into the regions $\{x>y\}$
and $\{x<y\}$ via the symmetry $k_L(x,y)=k_L(y,x)$.
This is consistent with the factor~$2$ in $J_{ij}(\tau) = 2\int_{-1}^{1-\tau}
P_i(x)P_j(x+\tau)\,dx$, which arises from the same kernel-symmetry splitting.
\end{remark}
The kernel quadratic form $K_{7/20}[v]$ reduces to
$2\int_0^2 (-\tfrac{7}{20})r''(\tfrac{7t}{20})\,C_v(t)\,dt$.
The integrand is split at $t = 10^{-4}$: the near-zero piece is bounded by
$|r''| < 2$ and Cauchy--Schwarz using exact rational arithmetic;
the remaining piece is computed by 256-bit Arb complex analytic integration
(\texttt{arb.integral}).

\begin{remark}[Verification standard for the integral bounds]
The auto-correlation polynomials and the near-zero piece are \emph{exact}
(pure rational arithmetic, machine-verified).
The remaining integral uses Arb~\cite{Johansson2017}, a peer-reviewed
C library for rigorous real and complex interval arithmetic; every Arb
result carries a certified enclosure radius.
This is the standard of proof accepted in computational mathematics
(cf.\ \textit{Mathematics of Computation}) and is qualitatively distinct
from unverified floating-point computation.
The Lean~4 formalisation (available at the supporting repository)
covers the integer comparisons via \texttt{native\_decide} and
additionally verifies $\log 2 < 7/10$ and $\sqrt{2} > 7/5$ using
Mathlib's \texttt{Real.sum\_le\_exp\_of\_nonneg} and
\texttt{Real.sqrt\_lt\_sqrt}. Formalising the Arb-based integral
in Lean~4 is a planned extension (Extended Goal~E3).
\end{remark}
Both upper endpoints of the resulting interval enclosures are strictly
below zero.
\end{proof}

%% ══════════════════════════════════════════════════════════════════════════════
\section{The Main Conjecture and Remaining Obligations}
\label{sec:conjecture}
%% ══════════════════════════════════════════════════════════════════════════════

\begin{definition}
\label{def:lambda}
For $L > 0$, let
\[
  \lambda(L) = \inf_{\substack{v \in D(Q_W^L) \\ v \ne 0}}
               \frac{Q_W^L(v)}{\norm{v}^2}.
\]
\end{definition}

\begin{remark}[Main conjecture FP-0.35]
It is \emph{conjectured} (but not proved in this paper) that
$\lambda(7/20) > 0$.
Because $\lambda$ is non-increasing in $L$, this would give
$\lambda(L) > 0$ for all $0 < L \le 7/20$.

The conjecture reduces, via \cref{cor:redistribution} and
\cref{thm:schur} with $N_+ = 8$, $N_- = 6$, $\eta = 1/2$, to the
verification of the two interval-matrix conditions:
\begin{align*}
  b_L > 0 \quad\text{and}\quad b_L F_{\rm even} - R_\eta^{\rm even} &\succ 0,\\
  b_L > 0 \quad\text{and}\quad b_L F_{\rm odd}  - R_\eta^{\rm odd}  &\succ 0.
\end{align*}
Discovery pilots (mpmath, dps=100) with $c_L = 0$ give
$\min\text{-pivot} \approx 0.529$ (even) and $\approx 0.560$ (odd),
but a formally certified proof requires closing the Archimedean
primitive chain (O2) and independently replaying the certificate.
\end{remark}

%% ══════════════════════════════════════════════════════════════════════════════
\section*{Acknowledgements}
The author thanks the open-source communities behind python-flint (Arb),
Lean~4 / Mathlib, and proofctl for the computational infrastructure
that underpins this work.
%% ══════════════════════════════════════════════════════════════════════════════

%% ══════════════════════════════════════════════════════════════════════════════
\begin{thebibliography}{99}

\bibitem{Johansson2017}
F.~Johansson, \emph{Arb: efficient arbitrary-precision midpoint-radius interval
arithmetic}, IEEE Trans.\ Comput.\ \textbf{66} (2017), no.~8, 1281--1292.

\bibitem{Bombieri2000}
E.~Bombieri, \emph{Remarks on Weil's quadratic functional in the theory of
prime numbers, I}, Rend. Mat. Acc. Lincei \textbf{11} (2000), 183--233.

\bibitem{Groskin2026}
A.~Groskin, \emph{A finite Guinand--Weil dictionary and archimedean tail order
for the truncated Weil quadratic form}, arXiv:2607.02828 (2026).

\bibitem{Kim2026}
T.~Kim et al., \emph{A numerical realization of Suzuki's Weil-quadratic-form
operator}, arXiv:2607.24830 (2026).

\bibitem{Suzuki2026}
M.~Suzuki, \emph{Weil's quadratic form via the screw function},
arXiv:2606.09096 (2026).

\bibitem{Weil1952}
A.~Weil, \emph{Sur les ``formules explicites'' de la th\'eorie des nombres
premiers}, Comm. S\'em. Math. Univ. Lund (1952), 252--265.

\bibitem{Yoshida1992}
M.~Yoshida, \emph{On the positivity of the Weil functional}, unpublished
manuscript (1992). The small-support positivity result is also discussed
in Bombieri~\cite{Bombieri2000}.

\end{thebibliography}

\end{document}
